\begin{frame}[shrink]
	\frametitle{O que é um peso, e como ele é calculado}
	\small
	\begin{itemize}
		\item Um \textbf{peso} é um número real associado a uma conexão entre dois neurônios --- é o parâmetro que a rede efetivamente \textbf{aprende}. Os dados de entrada não mudam; os pesos, sim.\newline
		
		\item Dentro de um neurônio, os pesos entram na \textbf{combinação linear}: cada entrada é multiplicada pelo seu peso, e os produtos são somados.
	\end{itemize}
	
	\begin{equation*}
		z = \sum_i w_i \, x_i \qquad y = \sigma(z)
	\end{equation*}
	
	\par Exemplo numérico real com o primeiro neurônio da rede do slide anterior:\newline
	\par Entrada $x = [1{,}0,\ 0{,}5,\ -0{,}5]$
	\par Pesos $w = [0{,}4,\ -0{,}3,\ 0{,}6]$.
	
	\begin{equation*}
		z = 0{,}4 \!\cdot\! 1{,}0 + (-0{,}3) \!\cdot\! 0{,}5 + 0{,}6 \!\cdot\! (-0{,}5) = -0{,}05
		\quad\Longrightarrow\quad
		y = \sigma(-0{,}05) \approx 0{,}4875
	\end{equation*}
	
	\par Cada um dos outros neurônios repete exatamente esse cálculo, com sua própria linha de pesos. Antes do treino, os pesos partem de algum valor inicial (aqui, fixos); o que o treino faz é ajustá-los, um pouco a cada passo, para que a saída da rede se aproxime do que se quer.
\end{frame}
